\appendix
\section{Методика тестирования и безопасность}

\subsection{Протокол безопасности полетов}

\subsubsection{1. Подготовка зоны (Pre-flight)}
\begin{itemize}
    \item Очистить зону $3 \times 3$ метра.
    \item Закрыть окна (ветер).
    \item Убрать животных и детей.
    \item Предупредить окружающих о запуске.
    \item Надеть защитные очки (для оператора).
\end{itemize}

\subsubsection{2. Проверка дрона (Checklist)}
\begin{itemize}
    \item \textbf{Пропеллеры:} Целые, затянуты, вращаются свободно.
    \item \textbf{Аккумулятор:} Заряжен > 3.8В на банку, закреплен надежно.
    \item \textbf{Камера:} Линза чистая, шлейф не болтается.
    \item \textbf{Wi-Fi:} Сигнал стабильный (RSS > -60 dBm).
\end{itemize}

\subsubsection{3. Экстренная ситуация}
Если дрон ведет себя неадекватно:
\begin{enumerate}
    \item Громко крикнуть <<ПОСАДКА!>>.
    \item Нажать кнопку LAND (если есть).
    \item Если нет реакции — аккуратно поймать снизу (за батарею) и перевернуть.
    \item Отключить аккумулятор.
\end{enumerate}

\subsection{Тестирование в разных условиях}

\subsubsection{Погодные условия (Ветер)}
Дрон Pioneer очень легкий (около 100 г). Он не предназначен для полетов на улице при ветре > 3 м/с.
Для теста в помещении используйте вентилятор:
\begin{itemize}
    \item Включите вентилятор на минимальную скорость с расстояния 2 метра.
    \item Попробуйте удержать точку (режим Optical Flow).
    \item Если дрон сносит > 50 см — условия непригодны для точной навигации.
\end{itemize}

\subsubsection{Освещение (Lux)}
Камера Pioneer имеет низкий динамический диапазон.
\begin{itemize}
    \item \textbf{Яркое солнце:} Пересвет (белое пятно). Маркеры не видны.
    \item \textbf{Сумерки:} Сильный шум матрицы. Маркеры <<дрожат>>.
    \item \textbf{Мерцание ламп:} Люминесцентные лампы (50 Гц) могут давать полосы на видео.
\end{itemize}

\section{Готовые конфигурации (Templates)}

\subsection{Файл конфигурации camera\_params.yaml}
\begin{codefiletitle}[bash]{camera_params.yaml}{python/examples/advanced_navigation/camera_params_yaml.yaml}
\end{codefiletitle}

\subsection{Шаблон скрипта полета (main.py)}
\begin{codefiletitle}[python]{Базовый шаблон}{python/examples/advanced_navigation/main_template_py.py}
\end{codefiletitle}
